\section{Established tests from galaxies to gravitational waves}
\label{sec:benchmarks}

The local acceleration relation, lensing spectra, and tensor propagation
test whether the effective branches recover measured phenomena already
described by MOND or conventional gravitating matter and general relativity.
Their selections and covariance treatments remain separate. Overlapping
catalogues and spectra require a joint likelihood before their evidence
can be combined. \source{U1}

\subsection{Local RAR and baryonic Tully--Fisher}
\label{sec:localrar}
For 130 selected SPARC galaxies with 2710 radial points, profiling the
stated stellar mass-to-light, distance, and inclination priors gives a
common-$a_0$ penalized score of 11258.0. Adding a
$5\,\kms$ velocity floor gives 3433.2, with median per-galaxy
$\chi^2$ per point of 1.54. The sensitivity to the adopted velocity
errors is substantial. This is a diagonal-error algebraic-RAR check;
it does not replace the full director-field calculation or account for
all correlated noncircular motions. For 100 flat-velocity galaxies,
the baryonic Tully--Fisher log residual has median 0.020 dex and rms
0.067 dex. \source{U1, S03}\cite{sparcnew}

\subsection{Galaxy weak lensing}
\label{sec:galaxylensing}
The isolated-lens audit distinguishes the older KiDS release with its
covariance from the updated published table. At the reference $a_0$,
the latter gives a diagonal statistic of 10.06 for fifteen bins, or
5.99 for seven inner bins, including the tabulated statistical and
interpolation/extrapolation uncertainties. Its full table covariance
was unavailable, and its published $H_0=73$ convention was retained.
The two reductions overlap and provide no independent combined evidence.
\source{U1, S04}\cite{kidsnew}


\subsection{Primary CMB spectra and projected Weyl power}
\label{sec:cmb}
The positive-allocation cosmology of Section~\ref{sec:cosmology} was
also evaluated against 135 ACT primary-spectrum bins, profiling the
stated calibration at fixed cosmological parameters. The score is
162.73 against 162.02 for the reference \LCDM\ fit. The ten-bin ACT
lensing comparison gives 16.21 against 17.86. A higher-accuracy
lensing calculation changes the bandpowers by 0.013\%. These are
separate checks; adding their scores to Planck would require the
appropriate overlap cuts and joint covariance. \source{U1, S05--S06}\cite{actnew}


The ACT lensing calculation tests projected Weyl power. A DESI full-shape
clustering likelihood with survey windows, bias, and redshift-space
distortions, and a nonlinear weak-lensing likelihood, remain open.
\source{U1, S06}

\subsection{Gravitational slip and the lensing--velocity statistic}
\label{sec:eg}
The lensing--velocity statistic $E_G$ depends on the Weyl potential and
the baryon velocity response. The calculation gives 0.4375 at
$z=0.316$ and 0.3862 at $z=0.555$, nearly constant over
$k=0.005$--$0.15\,\mathrm{Mpc}^{-1}$. The cited ACT--BOSS estimates
are $0.49^{+0.14}_{-0.11}$ and $0.31^{+0.06}_{-0.05}$.
The no-slip identity alone does not determine this observable.
In the conserved-fluid implementation, CAMB's ordinary
\texttt{delta\_tot} omits part of the bookkeeping capacity component;
a physical total-matter $\sigma_8$ must be reconstructed independently
of the arbitrary reference-dust split. \source{U1, S09}\cite{egnew}


\subsection{Tensor polarization and propagation}
\label{sec:tensor}
For a constant Einstein coefficient and minimally coupled scalar matter,
the tensor quadratic action has two positive-energy luminal modes.
Public GWTC-4 polarization comparisons give combined
$\log_{10}\mathcal B_{S/T}=-14.72$ and
$\log_{10}\mathcal B_{V/T}=-5.33$ for pure scalar and vector
alternatives. Mixed tensor--scalar and tensor--vector alternatives
have values $-0.20$ and $+0.10$, respectively. This supports
compatibility with the minimal tensor branch and leaves mixed sectors
weakly constrained; it supplies no distinctive preference over GR.
The native tensor continuum and a map from higher-curvature operators
to propagation or Lorentz-violation bounds remain open.
\source{U1, S12; D1}\cite{gwtcnew}

